\title{Group Sequence Policy Optimization}

\begin{abstract}
We present Group Sequence Policy Optimization (GSPO), a reinforcement learning method designed for the robust and effective training of large language models. GSPO differs from prior approaches that utilize token-level importance ratios by instead establishing sequence-level importance ratios, applying clipping, rewards, and optimization at the sequence level. Experimental results show that GSPO surpasses the GRPO algorithm in both efficiency and performance, offers substantial stabilization during Mixture-of-Experts (MoE) RL training, and holds promise for reducing complexity in RL infrastructure design. These advantages have played a key role in advancing the capabilities of the recent Qwen3 models.
\end{abstract}

\section{Introduction}
\label{sec:intro}

Reinforcement learning (RL) has become a central approach for advancing the performance of language models \citep{o1,dpsk-r1,qwq32b,qwen3}. By leveraging RL at scale, these models acquire the ability to address complex tasks—including competition-level mathematics and programming—by engaging in more thorough and extended reasoning.

Scaling RL effectively with increased computational resources requires that training remains both stable and resilient. Nevertheless, leading RL methods such as GRPO \citep{grpo} suffer from notable stability challenges during the training of very large language models, frequently leading to disastrous and irreversible collapses of the model \citep{qwen3, minimax-m1}. Such instability poses a major obstacle to further enhancing the capacities of language models through continued RL optimization.

This study reveals that GRPO's instability originates from a fundamental error in the application and legitimacy of importance sampling weights within its architecture. Such misuse gives rise to high-variance training noise, which escalates as the generated responses grow in length and is further intensified by the use of a clipping strategy, eventually causing the model to break down.

To overcome these foundational challenges, we introduce \textbf{Group Sequence Policy Optimization (GSPO)}, an RL algorithm tailored for training large language models. GSPO's central contribution is its principled formulation of the importance ratio derived from sequence likelihood \citep{click}, in strict accordance with the principles of importance sampling. Beyond that, GSPO calculates normalized rewards as the advantages among several generated responses for a given query, thereby reinforcing the coherence between reward assignment at the sequence level and the optimization process.

The results from our experiments highlight that GSPO substantially outperforms GRPO with respect to training stability, efficiency, and overall effectiveness. Notably, GSPO naturally addresses the common instability issues encountered during RL training of large-scale Mixture-of-Experts (MoE) architectures, thereby obviating the requirement for intricate stabilization measures, and suggests a pathway toward streamlining RL system design. These advantages have played a pivotal role in achieving remarkable enhancements within recent Qwen3 models. Looking forward, we regard GSPO as a dependable and scalable approach that can drive further progress in RL-based training methodologies for extensive language models.

\section{Preliminaries}

\paragraph{Notation}
Throughout this work, we represent the autoregressive language model with parameters $\theta$ as a policy $\pi_\theta$. The symbol $x$ refers to a query, and $\mathcal{D}$ indicates the collection of queries. For any response $y$ to a query $x$, the probability assigned by $\pi_\theta$ to generating $y$ given $x$ is expressed as $\pi_\theta (y | x)=\prod_{t=1}^{|y|} \pi_\theta (y_t | x, y_{<t} )$, where $|y|$ is the total number of tokens in $y$. The verifier $r$ can evaluate a query-response pair $(x, y)$, assigning it a reward $r(x, y) \in [ 0, 1 ]$.

\paragraph{Proximal Policy Optimization (PPO)} Using samples generated from the old policy $\pi_{\theta_\text{old}}$, PPO \citep{ppo} constrains the policy update within a proximal region of the old policy through the clipping mechanism. Specifically, PPO employs the following objective for policy optimization (we omit the KL regularization term hereinafter for brevity, as it is not the focus of this paper):

\begin{align}
\mathcal{J}_\text{PPO}(\theta) = \mathbb{E}_{ x \sim \mathcal{D},\, y \sim \pi_{\theta_\text{old}}( \cdot | x) }
\left[ \frac{1}{|y|} \sum_{t=1}^{|y|} 
\min \left( w_{t}(\theta) \widehat{A}_{t},  \, \mathrm{clip} \left( w_{t}(\theta), 1 - {\varepsilon}, 1 + {\varepsilon}\right) \widehat{A}_{t} \right)
\right],
\end{align}

In this context, the importance ratio for token $y_t$ is given by $
w_{t}(\theta) = \frac{ \pi_{\theta} (y_{t} | x, y_{<t}) }{ \pi_{\theta_\text{old}} (y_{t} | x,y_{<t})} $. The advantage $\widehat{A}_{t}$ associated with $y_t$ is obtained from a separate value model, and $\varepsilon$ specifies the range used for clipping these ratios.

A significant practical obstacle for PPO stems from its substantial dependence on the value model. Typically, the value model is comparable in scale to the policy model, resulting in increased demands on memory and computation. Moreover, the performance of PPO is closely tied to how accurately the value model estimates the advantage. Constructing a trustworthy value model is a difficult task by itself; scaling such a model to handle longer output sequences or more intricate assignments only amplifies this difficulty.

\paragraph{Group Relative Policy Optimization (GRPO)} GRPO \citep{grpo} bypasses the need for the value model by computing the relative advantage of each response within a group of responses to the same query. Specifically, GRPO optimizes the following objective:

\begin{align}
\label{equ:grpo}
\mathcal{J}_\text{GRPO}(\theta) = \mathbb{E}_{ x \sim \mathcal{D},\, \{y_i\}_{i=1}^G \sim \pi_{\theta_\text{old}}( \cdot | x) }
\left[ \frac{1}{G} \sum_{i=1}^{G} \frac{1}{|y_i|} \sum_{t=1}^{|y_i|} 
\min \left( w_{i,t}(\theta) \widehat{A}_{i,t},  \, \mathrm{clip} \left( w_{i,t}(\theta), 1 - {\varepsilon}, 1 + {\varepsilon}\right) \widehat{A}_{i,t} \right)
\right],
\end{align}

where $G$ is the number of generated responses to each query $x$ (i.e., the group size), and the importance ratio $w_{i,t}(\theta)$ and advantage $\widehat{A}_{i,t}$ of token $y_{i,t}$ are:

\begin{align}
    w_{i,t}(\theta)=\frac{ \pi_{\theta} (y_{i,t} | x, y_{i,<t}) }{ \pi_{\theta_\text{old}} (y_{i,t} | x,y_{i,<t})},\quad\ 
    \widehat{A}_{i,t} = \widehat{A}_{i} = \frac{r(x, y_i) - \mathrm{mean} \left( \{ r(x, y_i) \}_{i=1}^G \right) }{ \mathrm{std} \left( \{ r(x, y_i) \}_{i=1}^G \right) },
\end{align}

respectively, where all the tokens in $y_i$ share the same advantage as $\widehat{A}_{i}$.

\section{Motivation}
\label{sec:motivation}

The increasing scale of models, the prevalence of sparsity (as observed in Mixture-of-Experts models), and the extension of response lengths collectively drive the need for large rollout batch sizes in RL to fully leverage computational resources. To enhance the efficiency with which data is used, it is common to divide the extensive rollout dataset into several smaller mini-batches for each gradient computation. This practice inherently creates an off-policy learning environment, wherein responses $y$ originate from a previous policy $\pi_{\theta_\text{old}}$, rather than the present policy $\pi_\theta$ undergoing training. Consequently, this justifies the implementation of the clipping technique in PPO and GRPO, which serves to exclude excessively "off-policy" samples from contributing to the gradient calculation.

While mechanisms like clipping aim to manage this off-policy discrepancy, we identify a more fundamental issue in GRPO: \textit{its objective is ill-posed}. This problem becomes particularly acute when training large models on long-response tasks, leading to catastrophic model collapse. The ill-posed nature of the GRPO objective stems from a misapplication of importance sampling weights. The principle of importance sampling is to estimate the expectation of a function $f$ under a target distribution $\pi_\text{tar}$ by re-weighting samples drawn from a behavior distribution $\pi_\text{beh}$:

\begin{align}
\mathbb{E}_{z \sim \pi_\text{tar}} \left[ f(z) \right] = \mathbb{E}_{z \sim \pi_\text{beh}} \left[ \frac{ \pi_\text{tar} (z) }{ \pi_\text{beh} (z)} \, f(z) \right].
\end{align}

A central aspect of this method is the necessity to average across numerous samples ($N \gg 1$) drawn from the behavior distribution $\pi_\text{beh}$; this is essential for the importance weight $\frac{ \pi_\text{tar} (z) }{ \pi_\text{beh} (z)}$ to adequately rectify the mismatch between distributions.

By contrast, GRPO incorporates the importance weight $\frac{ \pi_{\theta} (y_{i,t} | x, y_{i,<t}) }{ \pi_{\theta_\text{old}} (y_{i,t} | x,y_{i,<t})}$ at every token position $t$. Here, the weight calculation is based solely on a single sample $y_{i,t}$ from each next-token distribution $\pi_{\theta_\text{old}}( \cdot | x, y_{i, <t})$. This approach is insufficient for effective distributional correction and instead generates substantial high-variance noise within the training gradients. Such variance compounds along longer sequences and is made worse by the use of a clipping mechanism. Empirical results indicate that these dynamics can cause irreversible model collapse. Should such collapse take place, recovery is typically unattainable—efforts such as restoring older checkpoints, fine-tuning hyperparameters (like clipping thresholds), extending output sequence lengths, or altering RL queries have not succeeded in reversing the failure.

This empirical outcome highlights a deeper flaw in the design of GRPO. The inadequacy of token-level importance weighting emphasizes an essential principle: \textbf{the granularity of the optimization objective must align with that of the reward signal}. Since the reward targets the complete sequence, it is evidently ineffective to apply off-policy correction at the token level. Thus, these findings direct us away from the token-level objective, instead motivating the direct use of importance weights and optimization at the \textit{sequence level}.

\section{Algorithm}

\subsection{GSPO: Group Sequence Policy Optimization}

Although the token-level importance weight $\frac{ \pi_{\theta} (y_{i,t} | x, y_{i,<t}) }{ \pi_{\theta_\text{old}} (y_{i,t} | x,y_{i,<t})}$ presents challenges in GRPO, it is noteworthy that, in language generation tasks, the \textit{sequence-level} importance weight $\frac{ \pi_{\theta} (y | x) }{ \pi_{\theta_\text{old}} (y | x)}$ possesses a well-defined theoretical interpretation: it quantifies the degree of divergence between the response $y$, as produced by $\pi_{\theta_\text{old}} (\cdot | x)$, and the distribution $\pi_{\theta} (\cdot | x)$. This property is consistent with the objectives of sequence-level rewards and serves as a relevant metric for evaluating the effects of the clipping mechanism.

Based on this straightforward observation, we propose the \textbf{Group Sequence Policy Optimization (GSPO)} algorithm. GSPO employs the following sequence-level optimization objective:

\begin{align}
\mathcal{J}_\text{GSPO} (\theta) =
\mathbb{E}_{ x \sim \mathcal{D},\, \{y_i\}_{i=1}^G \sim \pi_{\theta_\text{old}}( \cdot | x) }
\left[ 
\frac{1}{G} \sum_{i=1}^{G}
\min \left( s_{i}(\theta)  \widehat{A}_{i},  \, \mathrm{clip} \left( s_{i}(\theta), 1 - {\varepsilon}, 1 + {\varepsilon} \right) \widehat{A}_{i} \right) 
\right],
\label{equ:gspo}
\end{align}

where we adopt the group-based advantage estimation:

\begin{align}
\widehat{A}_{i} = \frac{r(x, y_i) - \mathrm{mean} \left( \{ r(x, y_i) \}_{i=1}^G \right) }{ \mathrm{std} \left( \{ r(x, y_i) \}_{i=1}^G \right) },
\end{align}

and define the importance ratio $s_{i}(\theta)$ based on sequence likelihood \citep{click}:

\begin{align}
s_{i}(\theta) = \left( \frac{ \pi_{\theta} (y_i | x) }{ \pi_{\theta_\text{old}} (y_i | x)} \right)^{\frac{1}{|y_i|}}
=
\exp \left( \frac{1}{|y_i|} \sum_{t=1}^{|y_i|} \log \frac{ \pi_{\theta} (y_{i,t} | x, y_{i,<t}) }{ \pi_{\theta_\text{old}} (y_{i,t} | x,y_{i,<t})} \right).
\end{align}

Consequently, GSPO implements clipping at the response level rather than at the token level, ensuring that highly ``off-policy'' responses are eliminated from the gradient estimation process. This approach is consistent with sequence-level reward assignment and optimization strategies. Additionally, length normalization is incorporated in $s_{i}(\theta)$ to mitigate variance and maintain $s_{i}(\theta)$ within a standardized numerical interval. Without this normalization, even minor changes in token likelihoods could cause substantial swings in the sequence-level importance ratios, necessitating different clipping thresholds for responses of varying lengths. It is important to highlight that the clipping intervals in GSPO, as compared to prior algorithms such as GRPO, often differ significantly in scale because of the varying conceptualizations of importance ratios.

\subsection{Gradient Analysis}
\label{subsec:gradient}

We can derive the gradient of the GSPO objective as follows (clipping is omitted for brevity):

\begin{align}
\nabla_{\theta} \mathcal{J}_\text{GSPO} (\theta)
=&\ 
\nabla_{\theta} \mathbb{E}_{ x \sim \mathcal{D},\, \{y_i\}_{i=1}^G \sim \pi_{\theta_\text{old}}( \cdot | x) }
\left[ 
\frac{1}{G} \sum_{i=1}^{G} s_{i}(\theta) \widehat{A}_{i}
\right] \\
= &\ 
\mathbb{E}_{ x \sim \mathcal{D},\, \{y_i\}_{i=1}^G \sim \pi_{\theta_\text{old}}( \cdot | x) }
\left[ 
\frac{1}{G} \sum_{i=1}^{G}
s_{i}(\theta) \widehat{A}_{i}
\cdot \nabla_{\theta} \log s_{i}(\theta)
\right] \\
=&\ 
\mathbb{E}_{ x \sim \mathcal{D},\, \{y_i\}_{i=1}^G \sim \pi_{\theta_\text{old}}( \cdot | x) }
\left[ 
\frac{1}{G} \sum_{i=1}^{G} 
\left( \frac{ \pi_{\theta} (y_i | x) }{ \pi_{\theta_\text{old}} (y_i | x)} \right)^{\frac{1}{|y_i|}} \widehat{A}_{i} 
\cdot \frac{1}{|y_i|} \sum_{t=1}^{|y_i|} \nabla_{\theta} \log \pi_{\theta} (y_{i,t} | x, y_{i,<t}) 
\right].
\label{equ:gspo_grad}
\end{align}

For comparison, the gradient of the GRPO objective is as follows (note that $\widehat{A}_{i,t} = \widehat{A}_{i}$):

\begin{align}
\nabla_{\theta} \mathcal{J}_\text{GRPO} (\theta) 
=&\ 
\nabla_{\theta} \mathbb{E}_{ x \sim \mathcal{D},\, \{y_i\}_{i=1}^G \sim \pi_{\theta_\text{old}}( \cdot | x) }
\left[ \frac{1}{G} \sum_{i=1}^{G} \frac{1}{|y_i|} \sum_{t=1}^{|y_i|} 
w_{i,t}(\theta) \widehat{A}_{i,t}
\right] \\
=&\
\mathbb{E}_{ x \sim \mathcal{D},\, \{y_i\}_{i=1}^G \sim \pi_{\theta_\text{old}}( \cdot | x) }
\left[ \frac{1}{G} \sum_{i=1}^{G} \widehat{A}_{i} 
\cdot \frac{1}{|y_i|} \sum_{t=1}^{|y_i|} 
\frac{ \pi_{\theta} (y_{i,t} | x, y_{i,<t}) }{ \pi_{\theta_\text{old}} (y_{i,t} | x,y_{i,<t})} 
\nabla_{\theta} \log \pi_{\theta} (y_{i,t} | x, y_{i,<t})  
\right].
\end{align}

Consequently, the primary difference between GSPO and GRPO pertains to \textit{the manner in which gradient contributions from the log likelihoods of individual tokens are weighted}. GRPO assigns each token a distinct ``importance weight'' given by $\frac{ \pi_{\theta} (y_{i,t} | x, y_{i,<t}) }{ \pi_{\theta_\text{old}} (y_{i,t} | x,y_{i,<t})}$, which may take values ranging from $(0, 1+\varepsilon]$ when $\widehat{A}_{i} >0$, or from $[1-\varepsilon, +\infty)$ when $\widehat{A}_{i} < 0$. These variable weights are significant and, as training advances, their cumulative effect can introduce considerable unpredictability. By comparison, GSPO applies uniform weighting to every token within a sequence, thereby removing the source of instability introduced by GRPO's heterogeneous weighting.

\subsection{GSPO-token: A Token-level Objective Variant}

In scenarios like multi-turn RL, we may desire a finer-grained advantage adjustment than the sequence level. To this end, we introduce a token-level objective variant of GSPO, namely \textbf{GSPO-token}, to allow token-wise advantage customization:

\begin{align}
\mathcal{J}_\text{GSPO-token}(\theta)
=
\mathbb{E}_{ x \sim \mathcal{D},\, \{y_i\}_{i=1}^G \sim \pi_{\theta_\text{old}}( \cdot | x) }
\left[ 
\frac{1}{G} \sum_{i=1}^{G} \frac{1}{|y_i|} \sum_{t=1}^{|y_i|} 
\min \left( s_{i,t}(\theta) \widehat{A}_{i,t},  \, \mathrm{clip} \left( s_{i,t}(\theta), 1 - {\varepsilon}, 1 + {\varepsilon} \right) \widehat{A}_{i,t} \right) 
\right],
\label{equ:gspo-token}
\end{align}

where

\begin{align}
s_{i,t}(\theta) 
= \mathrm{sg} \left[ s_{i}(\theta) \right]  \cdot \frac{ \pi_{\theta} (y_{i,t} | x, y_{i,<t}) }{ \mathrm{sg} \left[ \pi_{\theta} (y_{i,t} | x, y_{i,<t}) \right] },
\end{align}

and $\mathrm{sg}[\cdot]$ denotes only taking the numerical value but stopping the gradient, corresponding to the \texttt{detach} operation in PyTorch. The gradient of GSPO-token can be derived as:

\begin{align}
\nabla_{\theta} \mathcal{J}_\text{GSPO-token}(\theta)
=&\ 
\nabla_{\theta} \mathbb{E}_{ x \sim \mathcal{D},\, \{y_i\}_{i=1}^G \sim \pi_{\theta_\text{old}}( \cdot | x) }
\left[ 
\frac{1}{G} \sum_{i=1}^{G}
\frac{1}{|y_i|} \sum_{t=1}^{|y_i|} 
s_{i,t}(\theta) \widehat{A}_{i,t}
\right] \\
=&\ 
\mathbb{E}_{ x \sim \mathcal{D},\, \{y_i\}_{i=1}^G \sim \pi_{\theta_\text{old}}( \cdot | x) }
\left[ 
\frac{1}{G} \sum_{i=1}^{G} s_i (\theta) \cdot \frac{1}{|y_i|} \sum_{t=1}^{|y_i|} 
\widehat{A}_{i,t} \frac{ \nabla_{\theta} \pi_{\theta} (y_{i,t} | x, y_{i,<t}) }{ \pi_{\theta} (y_{i,t} | x, y_{i,<t}) }
\right] \\
=&\ 
\mathbb{E}_{ x \sim \mathcal{D},\, \{y_i\}_{i=1}^G \sim \pi_{\theta_\text{old}}( \cdot | x) }
\left[ 
\frac{1}{G} \sum_{i=1}^{G}  \left( \frac{ \pi_{\theta} (y_i | x) }{ \pi_{\theta_\text{old}} (y_i | x)} \right)^{\frac{1}{|y_i|}}
\cdot \frac{1}{|y_i|} \sum_{t=1}^{|y_i|}
\widehat{A}_{i,t} \nabla_{\theta} \log \pi_{\theta} (y_{i,t} | x, y_{i,<t})  
\right].
\label{equ:gspo-token_grad}
\end{align}

Observe that the expression $\frac{ \pi_{\theta} (y_{i,t} | x, y_{i,<t}) }{ \mathrm{sg} \left[ \pi_{\theta} (y_{i,t} | x, y_{i,<t}) \right] }$ evaluates to 1, which implies that $s_{i,t}(\theta)$ is, in fact, numerically equivalent to $s_{i}(\theta)$. When examining Equation~\eqref{equ:gspo} alongside Equation~\eqref{equ:gspo-token}, as well as Equation~\eqref{equ:gspo_grad} with Equation~\eqref{equ:gspo-token_grad}, we find that GSPO-token and GSPO yield identical results with respect to their optimization objectives, clipping criteria, and theoretical gradients, provided the advantage estimates for every token in the sequence $y_i$ are set uniformly (that is, $\widehat{A}_{i,t} = \widehat{A}_{i}$ for all $t$). Nonetheless, GSPO-token provides greater versatility, enabling individual token-level adjustment of the advantages.

\section{Experiments and Discussion}

\subsection{Empirical Results}

We conduct experiments utilizing a cold-start configuration by fine-tuning the Qwen3-30B-A3B-Base model, presenting the resulting training reward trajectories along with the model's performance curves across several benchmarks: AIME'24 (mean Pass@1 over 32 samples), LiveCodeBench (202410-202502, mean Pass@1 over 8 samples), and CodeForces (Elo Rating). Throughout RL training, each rollout batch is subdivided into four mini-batches for the purpose of gradient computation and updates. For GSPO, the left and right clipping parameters in Equation~\eqref{equ:gspo} are designated as 3e-4 and 4e-4, correspondingly. GRPO is employed as the baseline comparator; its left and right clipping intervals for Equation~\eqref{equ:grpo} are set to 0.2 and 0.27, respectively, following extensive tuning to achieve an equitable comparison. It is important to highlight that successful convergence of MoE RL under GRPO requires the Routing Replay strategy, which is addressed further in \S~\ref{subsec:routing_replay}, whereas \textit{GSPO eliminates the dependency on this approach}.

As illustrated in Figure~\ref{fig:results}, GSPO enables consistently stable progress during training. Our results indicate that \textbf{GSPO facilitates sustained advancements in model performance by scaling up computational resources for training, systematically refreshing the query set, and increasing generation length}. In addition, GSPO exhibits markedly higher training efficiency compared to GRPO, surpassing it in both training accuracy and benchmark results when provided with equal training compute and identical numbers of queries. Lastly, GSPO has been successfully implemented for RL training with the latest Qwen3 models, offering compelling evidence of GSPO’s effectiveness in harnessing RL scaling for large language models.

\subsection{Curious Observation on Clipping Fractions}

One notable difference between GSPO and GRPO lies in how responses are clipped: GSPO applies clipping to whole sequences, whereas GRPO clips at the token level. Figure~\ref{fig:clipping} illustrates that the proportion of clipped tokens in GSPO versus GRPO differs by nearly two orders of magnitude; moreover, modifying the clipping intervals does not affect this significant discrepancy. Surprisingly, GSPO, which clips many more tokens—thereby reducing the number available for training or gradient estimation—still demonstrates greater training efficiency than GRPO. This seemingly paradoxical result, in which clipping a substantially higher fraction of tokens yields improved efficiency, suggests that the token-level gradients provided by GRPO are both noisier and less effective for extracting useful information from the data. By contrast, GSPO’s sequence-level method yields a more robust and informative training signal.

\subsection{Benefit of GSPO for MoE Training}
\label{subsec:routing_replay}

\paragraph{Background}
MoE models, due to their sparsely activated architecture, encounter distinct stability issues during RL training that are not present in dense models. Specifically, our observations indicate that utilizing the GRPO algorithm can exacerbate \textit{expert-activation volatility} within MoE models, thereby impeding the proper convergence of RL training. More concretely, even after a few gradient steps, the set of experts selected for the same response can change considerably. For instance, when training the 48-layer Qwen3-30B-A3B-Base model, we observed that, for the same rollout sample, approximately 10\% of the experts chosen under the updated policy $\pi_{\theta}$ differ from those used by the preceding policy $\pi_{\theta_\text{old}}$ after each RL update. This issue becomes increasingly pronounced in deeper MoE architectures and results in highly unstable token-level importance ratios, defined as $w_{i,t}(\theta) = \frac{ \pi_{\theta} (y_{i,t} | x, y_{i,<t}) }{ \pi_{\theta_\text{old}} (y_{i,t} | x,y_{i,<t})}$. The pronounced oscillations in these values undermine their reliability, as elaborated in \S~\ref{sec:motivation} and~\ref{subsec:gradient}, thereby obstructing the typical convergence behavior expected during RL training.

\paragraph{Our Previous Approach} In our earlier work, we addressed this problem by implementing the \textbf{Routing Replay} method during training. This involves storing the set of activated experts as determined by $\pi_{\theta_\text{old}}$, and subsequently ``replaying'' these expert activation routes in $\pi_{\theta}$ for the purpose of calculating importance ratios: $w_{i,t}(\theta) = \frac{ \pi_{\theta} (y_{i,t} | x, y_{i,<t}) }{ \pi_{\theta_\text{old}} (y_{i,t} | x,y_{i,<t})}$. Through this mechanism, for each output token $y_{i,t}$, both $\pi_{\theta} (y_{i,t} | x, y_{i,<t})$ and $\pi_{\theta_\text{old}} (y_{i,t} | x, y_{i,<t})$ utilize identical expert selection pathways, which promotes stable token-level importance ratios and facilitates the optimization of the same activated subnetwork across successive gradient steps. Figure~\ref{fig:routing_replay} illustrates that Routing Replay is a critical component for achieving stable and proper convergence when training Mixture of Experts models with GRPO.

\paragraph{Benefit of GSPO} While Routing Replay facilitates proper GRPO training and convergence for MoE models, its approach of recycling routing modes introduces supplementary memory consumption and communication demands, potentially restricting the effective capacity of the MoE architecture. On the other hand, GSPO, as depicted in Figure~\ref{fig:results}, completely removes reliance on Routing Replay. GSPO can directly compute importance ratios $s_i(\theta)$ and proceed with regular convergence and stable optimization. The core advantage lies in GSPO's exclusive reliance on sequence likelihood (i.e., $\pi_{\theta} (y_i | x)$), rendering it largely indifferent to individual token likelihoods (i.e., $\pi_{\theta} (y_{i,t} | x, y_{i,<t})$). Given that MoE models consistently preserve their ability for language modeling, sequence likelihood exhibits only moderate variation. Thus, GSPO addresses the instability of expert activation in MoE architectures at its root, dispensing with elaborate methods such as Routing Replay. This approach not only streamlines and fortifies training but also empowers the MoE model to realize its complete potential, free from imposed limitations.

\subsection{Benefit of GSPO for RL Infrastructure}

Due to inconsistencies in numerical precision between training systems (such as Megatron) and inference systems (such as SGLang and vLLM), it is common practice to calculate the likelihoods of generated responses under the old policy $\pi_{\theta_\text{old}}$ using the training system. Nevertheless, GSPO relies on sequence-level rather than token-level likelihoods during optimization. Conceptually, sequence-level likelihoods are more robust to precision differences. This characteristic allows GSPO to utilize likelihoods produced by inference systems directly for optimization, thereby eliminating the necessity of recomputing them with the training system. This approach proves advantageous, particularly in situations involving partial rollout, multi-turn reinforcement learning, and settings with separate training and inference infrastructures.

\section{Conclusion}

We introduce Group Sequence Policy Optimization (GSPO), a novel reinforcement learning method tailored for large language model training. GSPO adheres to the core tenets of importance sampling by establishing sequence-likelihood-based importance ratios, enabling sequence-level clipping, reward assignment, and optimization steps. Our results indicate that GSPO offers significant enhancements in training robustness, efficiency, and overall effectiveness relative to GRPO, especially excelling in large-scale reinforcement learning contexts for MoE models. This advancement has directly contributed to the marked gains observed in the recent Qwen3 models. With GSPO providing a robust and scalable foundation, we are poised to extend RL methods further and anticipate major progress in artificial intelligence stemming from this trajectory.

\end{document}